\documentclass[11pt]{article}
\input{style-pc}
\usepackage{array}

\begin{document}

\entetesujet{Sujet bac 2014 -- Physique-Chimie -- Série C}

% ============================ CHIMIE ============================
\matiere{CHIMIE}{8 points}

\partie{A : vérification des connaissances}

\noindent\textbf{I. Questions à choix multiples.} Choisis la bonne réponse.

\begin{enumerate}
  \item Le pH d'une solution de dibase forte, de concentration $C_b$ est :
  \begin{enumerate}[label=\alph*)]
    \item $-\log C_b$ ;
    \item $-\log 2C_b$ ;
    \item $14 + \log 2C_b$ ;
    \item $14 + \log C_b$.
  \end{enumerate}
  \item La réaction d'estérification est une réaction :
  \begin{enumerate}[label=\alph*)]
    \item exothermique ;
    \item athermique ;
    \item thermique ;
    \item endothermique.
  \end{enumerate}
  \item L'oxydation est une réaction chimique qui correspond à :
  \begin{enumerate}[label=\alph*)]
    \item la diminution du nombre d'oxydation ;
    \item l'augmentation du nombre d'oxydation.
  \end{enumerate}
  \item Lors d'une réaction d'hydrolyse, on utilise un catalyseur pour :
  \begin{enumerate}[label=\alph*)]
    \item ralentir la réaction ;
    \item arrêter la réaction ;
    \item accélérer la réaction ;
    \item modifier la composition de la réaction.
  \end{enumerate}
\end{enumerate}

\noindent\textbf{II. Réponds par vrai ou faux.}
\begin{enumerate}
  \item La désintégration $\alpha$ se produit avec des noyaux lourds.
  \item Lors de l'absorption, l'atome d'hydrogène passe d'un niveau supérieur vers un niveau inférieur.
  \item L'abaissement cryométrique est proportionnel à la concentration $C$ de la solution.
  \item L'énergie d'un atome dans son état fondamental est maximale.
\end{enumerate}

\noindent\textbf{III. Texte à trous.} Complète la phrase suivante en remplaçant les quatre mots manquants par les mots suivants : \textit{réaction ; équilibrée ; réactionnel ; coexistent}.

\textit{Une réaction $\cdots\cdots$ est une $\cdots\cdots$ au cours de laquelle les réactifs et les produits $\cdots\cdots$ dans le milieu $\cdots\cdots$}

\partie{B : application des connaissances}
Le radon $\ch{^{222}_{86}Rn}$ a une période ou demi-vie de $3{,}8$ jours. Il est radioactif $\alpha$.

\begin{enumerate}
  \item Écris l'équation bilan de sa désintégration.
  \item Calculer sa constante radioactive.
  \item On dispose d'un échantillon de $0{,}20$ mg de radon 222. Combien y a-t-il de noyaux radioactifs dans l'échantillon ?
  \item Quelle est l'activité de l'échantillon ?
  \item Quelle sera l'activité de l'échantillon au bout de 20 jours ?
\end{enumerate}
\noindent On donne $\aleph = 6{,}02\cdot10^{23}$ mol$^{-1}$ ; $M(\ch{Rn}) = 222$ g$\cdot$mol$^{-1}$.

\noindent Extrait du tableau périodique :
\begin{center}\renewcommand{\arraystretch}{1.3}
\begin{tabular}{|l|c|c|c|c|c|c|}
\hline
Nom & Bismuth & Polonium & Astate & Radon & Francium & Radium \\
\hline
Symbole & $\ch{_{83}Bi}$ & $\ch{_{84}Po}$ & $\ch{_{85}At}$ & $\ch{_{86}Rn}$ & $\ch{_{87}Fr}$ & $\ch{_{88}Ra}$ \\
\hline
\end{tabular}
\end{center}

% ============================ PHYSIQUE ============================
\matiere{PHYSIQUE}{12 points}

\partie{A : vérification des connaissances}

\noindent\textbf{1. Réponds par vrai ou faux.}
\begin{enumerate}[label=\alph*.]
  \item L'accélération du mouvement d'un objet en chute libre dépend de sa masse.
  \item Dans un pendule conique, l'angle d'écartement $\theta$ du fil par rapport à l'axe vertical est lié à sa vitesse angulaire $\omega$ par la relation : $\dfrac{1}{\cos\theta} = \dfrac{\omega^2 L}{g}$.
  \item L'équation différentielle d'un pendule élastique est de la forme :
  \begin{enumerate}[label=c\textsubscript{\arabic*})]
    \item $\dfrac{d^2x}{dt^2} + \dfrac{m}{k}x = 0$.
    \item $\dfrac{d^2x}{dt^2} + kmx = 0$.
    \item $\dfrac{d^2x}{dt^2} + \dfrac{k}{m}x = 0$.
  \end{enumerate}
\end{enumerate}

\noindent\textbf{2. Réarrangement (texte en désordre).} La phrase suivante concernant la définition de l'interfrange a été écrite en désordre. Ordonne-la.

\textit{est / de même / la distance / qui sépare / nature / l'interfrange / de deux / franges consécutives / les milieux.}

\partie{B : application des connaissances}

\noindent\textbf{Exercice (Effet photoélectrique).} Une cellule photoélectrique à cathode métallique (strontium) est éclairée simultanément par deux radiations monochromatiques de fréquences respectives $\nu_1 = 6{,}66\cdot10^{14}$ Hz et $\nu_2 = 4{,}84\cdot10^{14}$ Hz. Le seuil photoélectrique de la cathode est : $\lambda_0 = 0{,}6$ µm.

\begin{enumerate}
  \item Quelle est, de ces deux radiations, celle qui provoque l'effet photoélectrique ?
  \item Calcule la vitesse maximale avec laquelle un électron sort de la cathode.
  \item Le rendement quantique de la cellule étant $r = 0{,}02$ et l'intensité du courant de saturation $I_S = 10^{-6}$ A, déterminer la puissance lumineuse reçue par la cathode.
\end{enumerate}
\noindent On donne : $h = 6{,}62\cdot10^{-34}$ J$\cdot$s ; $C = 3\cdot10^8$ m$\cdot$s$^{-1}$ ; $m = 9{,}1\cdot10^{-31}$ kg (masse de l'électron) ; $e = 1{,}6\cdot10^{-19}$ C.

\partie{C : résolution d'un problème}
Un élève de terminale veut déterminer les coordonnées d'une bille $(B)$ au point de chute après rupture du fil de suspension. Pour cela, il dispose d'un pendule simple constitué d'un fil inextensible et sans masse de longueur $l = 2{,}0$ m, portant à son extrémité inférieure une petite bille $(B)$ de masse $m = 100$ g. La bille $(B)$ est considérée comme ponctuelle. L'autre extrémité du fil est fixée à un support en un point $O$. À l'équilibre, le pendule est vertical et la bille se trouve alors à une hauteur $h = 2{,}5$ m du sol. On prendra $g = 9{,}8$ m/s$^2$.

\begin{center}
\begin{tikzpicture}[>=Stealth,scale=1]
  \coordinate (O) at (0,0);
  \coordinate (A) at (0,-3);
  \coordinate (B) at ({3*sin(58)},{-3*cos(58)});
  \coordinate (sol) at (0,-4.3);
  % verticale dashed
  \draw[dashed] (O)--(0,-4.3);
  % fil incliné
  \draw[thick] (O)--(B);
  % angle
  \draw (0,-0.8) arc[start angle=270,end angle={270+58},radius=0.8];
  % axe x horizontal en A
  \draw[->,dashed] (3.6,-3)--(-3.6,-3);
  % sol
  \fill[pattern=north east lines] (-3.6,-4.45) rectangle (3.8,-4.3);
  \node[below right] at (2.6,-4.3){Sol};
  % points
  \fill (O) circle (1.2pt) node[above]{$O$};
  \fill (B) circle (2pt) node[right]{$(B)$};
  \fill (A) circle (1pt) node[above right]{$A$};
  % h
  \draw[<->] (-0.18,-3)--(-0.18,-4.3) node[midway,left]{$h$};
\end{tikzpicture}
\end{center}

On écarte le pendule d'un angle $\alpha_0 = 60^\circ$ de sa position d'équilibre et on le lâche sans vitesse initiale.

\begin{enumerate}
  \item \begin{enumerate}[label=\alph*.]
    \item Exprimer en fonction de $g$, $l$ et $\alpha_0$, la vitesse $v$ de la bille $(B)$ au passage par la position d'équilibre. En déduire sa valeur numérique.
    \item Déterminer l'intensité $T$ de la tension du fil lorsque le pendule passe par la verticale.
  \end{enumerate}
  \item Lorsque la bille passe par la verticale, le fil de suspension se coupe. La bille effectue un mouvement de chute et arrive au sol en un point $C$.
  \begin{enumerate}[label=\alph*.]
    \item Établis l'équation cartésienne de la trajectoire dans le repère $(A,\vi,\vj)$.
    \item Détermine les coordonnées $x_C$ et $y_C$ du point $C$, lieu de chute de la bille sur le sol.
  \end{enumerate}
\end{enumerate}

\end{document}
